\chapter{Stadiul actual al cercetării}

Acest capitol trece în revistă literatura de specialitate privind valorile lipsă în seriile de timp și utilizarea tehnicilor de învățare automată pentru imputare. Sunt analizate principalele direcții de cercetare, cu metodele propuse pentru tratarea datelor incomplete. Performanța modelelor de învățare automată depinde de calitatea datelor de intrare, de mecanismul prin care apar valorile lipsă și de modul în care acestea sunt tratate înainte de antrenare.~\cite{rubin1976inference, little2002missing, schafer2002missing}.

Capitolul acoperă atât contribuțiile clasice privind tratarea valorilor lipsă, cât și lucrările care folosesc rețele neuronale recurente pentru prognoza seriilor de timp. Această trecere în revistă situează lucrarea curentă în raport cu literatura existentă și arată unde se plasează comparația dintre modelele RNN, LSTM și GRU în condiții de date incomplete.

\section{Valori lipsă în seriile de timp}

Valorile lipsă apar frecvent în datele reale, inclusiv în seriile de timp. Cauzele pot fi variate: erori de măsurare, defectarea echipamentelor de achiziție, întreruperi în transmiterea datelor sau eliminarea unor observații considerate nevalide. În literatura de specialitate, problema nu este privită doar ca absența unor observații, ci și prin efectul pe care îl are asupra structurii temporale a datelor, asupra estimării parametrilor și asupra performanței predictive a modelelor.

Acock arată că modul de tratare a valorilor lipsă afectează direct validitatea concluziilor statistice și robustețea analizelor ulterioare~\cite{acock2005missing}. În cazul seriilor de timp, problema este mai dificilă decât în date tabulare obișnuite: absența unor observații nu perturbă doar distribuția variabilelor, ci și dependențele dintre momentele succesive ale seriei. De aceea, metodele de completare trebuie alese în funcție de mecanismul de apariție a valorilor lipsă și de natura secvențială a datelor.

Mecanismele de apariție a valorilor lipsă sunt de obicei clasificate în \textit{MCAR}, \textit{MAR} și \textit{MNAR}, iar strategia de tratare depinde de aceste ipoteze și de contextul de colectare~\cite{rubin1976inference, little2002missing}.

\section{Metode consacrate de tratare a valorilor lipsă}

Literatura clasică și aplicativă descrie numeroase metode de tratare a valorilor lipsă. Cele mai simple includ eliminarea observațiilor incomplete, completarea cu zero, completarea cu media sau mediana, propagarea ultimei valori observate și metodele bazate pe medii mobile sau ferestre glisante. Acestea sunt ușor de implementat și au cost computațional redus, dar pot introduce distorsiuni semnificative în structura seriei, mai ales când datele au o dinamică temporală complexă~\cite{little2002missing, schafer2002missing}.

Metodele simple funcționează rezonabil când procentul valorilor lipsă este mic sau când analiza este exploratorie. Devin insuficiente când absența datelor afectează consistent seria sau când dependențele temporale trebuie păstrate. Ca urmare, o bună parte a cercetărilor mai recente s-a îndreptat spre metode de imputare mai sofisticate, care folosesc relațiile dintre variabile și istoricul temporal pentru a reconstrui valorile absente.

\section{Metode moderne de imputare și abordări predictive}

Literatura propune tehnici de imputare bazate pe modele statistice iterative, algoritmi de tip ensemble și metode de învățare automată. Aici se încadrează imputarea multiplă, metodele bazate pe k-nearest neighbors și tehnicile predictive care folosesc arbori de decizie sau păduri aleatoare pentru a estima valorile lipsă. Toate exploatează structura datelor și corelațiile dintre variabile pentru a produce imputări mai precise.

Pentru seriile de timp, avantajul principal este că păstrează mai bine coerența datelor folosite ulterior la antrenare. Există și limitări: unele metode sunt costisitoare computațional, altele sunt sensibile la alegerea hiperparametrilor sau introduc bias când mecanismul valorilor lipsă nu este corect identificat. Compararea metodelor de imputare rămâne o direcție activă de cercetare, mai ales când scopul final este prognoza seriilor de timp~\cite{buuren2011mice, stekhoven2012missforest, breiman2001rf, che2018rnnmissing}.

\section{Învățarea federată în modelarea seriilor de timp}

Învățarea federată este o paradigmă în care modelele sunt antrenate pe date distribuite, iar actualizările locale sunt combinate fără a centraliza toate observațiile~\cite{mcmahan2017fedavg, kairouz2021fl}. Propusă inițial pentru scenarii distribuite de tip general, este utilă și pentru serii de timp, unde sursele de date sunt separate natural pe utilizatori, dispozitive, locații sau intervale temporale.

Principalul avantaj este compatibilitatea cu cerințele de confidențialitate și cu limitările de transfer al datelor. Dificultatea principală apare când distribuțiile locale diferă semnificativ, ceea ce poate afecta convergența și calitatea modelului global~\cite{kairouz2021fl}. În contextul forecasting-ului, problema este concretă: seriile distribuite între mai mulți utilizatori nu au mereu aceeași scară, volatilitate sau ritm temporal.

Relevanța pentru tema lucrării constă în faptul că oferă un cadru comparativ în care modelele locale sunt antrenate separat, iar performanța poate fi evaluată atât la nivel local, cât și global. Literatura dedicată forecasting-ului pe serii de timp distribuite rămâne mai restrânsă față de cea generală despre federated learning, însă numărul studiilor în această direcție crește, mai ales în aplicații unde datele sunt fragmentate între surse multiple.

\section{Ansamble și agregare de modele în forecasting}

Metodele de ansamblu sunt frecvent utilizate în literatura de predicție pentru că reduc dependența de un singur model și pot stabiliza rezultatele în prezența variației din date~\cite{dietterich2000ensemble, zhou2012ensemble}. În forma clasică, agregă mai multe predicții, iar performanța finală depinde de diversitatea membrilor ansamblului și de modul în care sunt combinați.

În cazul seriilor de timp, ansamblele sunt folosite atât pentru prognoză directă, cât și pentru sisteme hibride în care mai multe modele secvențiale sunt combinate sau reziduurile unui model de bază sunt corectate. Câștigul este mai mare când o singură arhitectură nu surprinde complet nelinearitățile sau când intrarea conține zgomot și variație locală ridicată.

Un concept central în literatura despre ansamble este diversitatea: modelele individuale trebuie să greșească diferit pentru ca agregarea să aducă un avantaj real~\cite{dietterich2000ensemble}. Aceasta este premisa de la care pornește și experimentul de ansamblu din lucrarea de față, care compară strategii de combinare și evaluează impactul lor asupra metricilor de predicție.

\section{Rețele recurente în predicția seriilor de timp}

Rețelele neuronale recurente au ajuns să fie larg folosite în modelarea datelor secvențiale odată cu maturizarea tehnicilor de învățare profundă. RNN, LSTM și GRU sunt arhitecturi frecvent aplicate în predicția seriilor de timp, datorită capacității de a surprinde dependențele temporale și de a integra informația istorică în predicție. Rezultatele din literatură arată că depășesc metodele tradiționale mai ales când seria prezintă nelinearități, variații locale puternice sau structuri temporale complexe.

Fiecare arhitectură vine cu avantaje și limitări distincte. RNN standard are o structură simplă, dar are dificultăți în captarea dependențelor pe termen lung. LSTM introduce mecanisme explicite de memorie și porți interne, ceea ce îl face mai robust pe secvențe lungi, dar crește complexitatea modelului. GRU păstrează avantajele modelării recurente cu o arhitectură mai simplă și este adesea preferat în aplicații practice unde eficiența computațională contează~\cite{sherstinsky2020rnn, yu2019lstmreview, korstanje2025gru, cho2014gru}.

\section{Studii comparative și limitări identificate în literatură}

Există numeroase studii care compară arhitecturile recurente din perspectiva acurateței predictive, a vitezei de antrenare și a capacității de generalizare. Rezultatele nu indică un model universal superior, deoarece performanța depinde de domeniu, de dimensiunea datelor, de strategia de preprocesare și de calitatea seriilor. Multe lucrări se concentrează pe alegerea arhitecturii sau pe reglarea hiperparametrilor, tratând problema valorilor lipsă sumar, de regulă printr-o singură metodă de completare aplicată înainte de antrenare.

Aceasta este una dintre limitările importante ale stadiului actual al cercetării. Există contribuții separate privind imputarea datelor și contribuții separate privind predicția seriilor de timp cu rețele recurente, dar puține studii analizează sistematic relația dintre metoda de tratare a valorilor lipsă și performanța finală a modelelor. Literatura oferă soluții pe fiecare parte, dar mai rar le pune față în față: efectul concret al diferitelor strategii de imputare asupra arhitecturilor RNN, LSTM și GRU în același cadru experimental rămâne insuficient studiat.

\section{Poziționarea lucrării curente}

Față de literatura analizată, lucrarea de față urmărește să evalueze concret cum diferite strategii de tratare a valorilor lipsă influențează performanța arhitecturilor RNN, LSTM și GRU. Abordarea este comparativă și aplicativă: aceleași date, aceleași condiții de antrenare, strategii de imputare diferite, arhitecturi diferite. Lucrarea se situează la joncțiunea dintre analiza datelor incomplete și modelarea predictivă a seriilor de timp prin metode de învățare automată.

Cadrul experimental unitar, în care datele incomplete, metodele de imputare și arhitecturile recurente sunt analizate împreună, este contribuția principală față de lucrările existente. Utilizarea unor seturi de date cu caracteristici diferite permite verificarea dacă concluziile se mențin dincolo de un singur context aplicativ. Capitolele următoare detaliază metodologia și experimentele care decurg din această perspectivă.
